% Options for packages loaded elsewhere
\PassOptionsToPackage{unicode}{hyperref}
\PassOptionsToPackage{hyphens}{url}
\PassOptionsToPackage{dvipsnames,svgnames,x11names}{xcolor}
\documentclass[
  11pt,
  a4paper,
]{article}
\usepackage{xcolor}
\usepackage[margin=2.6cm]{geometry}
\usepackage{amsmath,amssymb}
\setcounter{secnumdepth}{5}
\usepackage{iftex}
\ifPDFTeX
  \usepackage[T1]{fontenc}
  \usepackage[utf8]{inputenc}
  \usepackage{textcomp} % provide euro and other symbols
\else % if luatex or xetex
  \usepackage{unicode-math} % this also loads fontspec
  \defaultfontfeatures{Scale=MatchLowercase}
  \defaultfontfeatures[\rmfamily]{Ligatures=TeX,Scale=1}
\fi
\usepackage{lmodern}
\ifPDFTeX\else
  % xetex/luatex font selection
  \setmainfont[]{Libertinus Serif}
  \setmonofont[Scale=0.85]{Latin Modern Mono}
  \setmathfont[]{Latin Modern Math}
\fi
% Use upquote if available, for straight quotes in verbatim environments
\IfFileExists{upquote.sty}{\usepackage{upquote}}{}
\IfFileExists{microtype.sty}{% use microtype if available
  \usepackage[]{microtype}
  \UseMicrotypeSet[protrusion]{basicmath} % disable protrusion for tt fonts
}{}
\makeatletter
\@ifundefined{KOMAClassName}{% if non-KOMA class
  \IfFileExists{parskip.sty}{%
    \usepackage{parskip}
  }{% else
    \setlength{\parindent}{0pt}
    \setlength{\parskip}{6pt plus 2pt minus 1pt}}
}{% if KOMA class
  \KOMAoptions{parskip=half}}
\makeatother
\usepackage{longtable,booktabs,array}
\newcounter{none} % for unnumbered tables
\usepackage{calc} % for calculating minipage widths
% Correct order of tables after \paragraph or \subparagraph
\usepackage{etoolbox}
\makeatletter
\patchcmd\longtable{\par}{\if@noskipsec\mbox{}\fi\par}{}{}
\makeatother
% Allow footnotes in longtable head/foot
\IfFileExists{footnotehyper.sty}{\usepackage{footnotehyper}}{\usepackage{footnote}}
\makesavenoteenv{longtable}
\setlength{\emergencystretch}{3em} % prevent overfull lines
\providecommand{\tightlist}{%
  \setlength{\itemsep}{0pt}\setlength{\parskip}{0pt}}

\providecommand{\E}{\mathbb{E}}
\providecommand{\N}{\mathbb{N}}
\providecommand{\Z}{\mathbb{Z}}
\providecommand{\R}{\mathbb{R}}
\providecommand{\eps}{\varepsilon}
\providecommand{\ceil}[1]{\left\lceil #1\right\rceil}
\providecommand{\floor}[1]{\left\lfloor #1\right\rfloor}
\AtBeginDocument{\let\setminus\smallsetminus}
\usepackage[most]{tcolorbox}
\newtcolorbox{warning}{colback=red!5,colframe=red!65!black,boxrule=0.8pt,arc=2pt,left=6pt,right=6pt,top=5pt,bottom=5pt}
\usepackage{etoolbox}
\AtBeginEnvironment{longtable}{\small}
\setlength{\emergencystretch}{3em}
\usepackage{fancyhdr}
\pagestyle{fancy}
\fancyhf{}
\fancyhead[L]{\small\itshape Candidate result 2603.02786\_\_00 --- GPT-6 Astra, awaiting review}
\fancyhead[R]{\small\thepage}
\renewcommand{\headrulewidth}{0.2pt}
\usepackage{bookmark}
\IfFileExists{xurl.sty}{\usepackage{xurl}}{} % add URL line breaks if available
\urlstyle{same}
\hypersetup{
  pdftitle={A candidate proof of Conjecture 1 in Packing arithmetic progressions},
  pdfauthor={github.com/graph-theory-AI --- machine-generated candidate writeup},
  colorlinks=true,
  linkcolor={blue!50!black},
  filecolor={Maroon},
  citecolor={Blue},
  urlcolor={blue!50!black},
  pdfcreator={LaTeX via pandoc}}

\title{A candidate proof of Conjecture 1 in Packing arithmetic
progressions}
\author{\href{https://github.com/graph-theory-AI}{github.com/graph-theory-AI}
--- machine-generated candidate writeup}
\date{16 September 2026}

\begin{document}
\maketitle
\begin{abstract}
Fixed-size CRT-generic prime blocks, together with clustered residue
classes for their multiples and a tail decomposition, give m(n)
\textasciitilde{} (4/3)n\^{}(3/2)/log n.~This candidate proof was
generated by GPT-6 Astra in a single model pass for the Graph Theory AI
project. It was selected from the campaign because the model marked it
\texttt{would\_publish:\ true}; it has not been checked by the
adversarial referee, a human mathematician, or a novelty search.
\end{abstract}

\section{The problem and the claimed
result}\label{the-problem-and-the-claimed-result}

This document concerns
\href{https://graph-theory-ai.github.io/graph-conjectures/arxiv/2603.02786__00/}{Conjecture
1} from \emph{Packing arithmetic progressions}.

\subsection{Problem statement}\label{problem-statement}

Conjecture. \(m(n)\sim\frac{4}{3}\cdot\frac{n^{3/2}}{\ln n}\).

\subsection{Catalog context}\label{catalog-context}

The paper proves \(m(n)=\Theta(n^{3/2}/\ln n)\) with matching constants
\(4/3\) (lower) and \(5/3\) (upper) in Theorem 1. The authors state that
closing this gap and pinning down the exact constant requires additional
ideas, and suspect \(4/3\) is correct.

Also stated in

Packing arithmetic progressions (2026-03-03)

Packing arithmetic progressions (2026-03-03)

\subsection{Claimed result}\label{claimed-result}

\textbf{Model verdict: proved} (confidence: high).

Fixed-size CRT-generic prime blocks, together with clustered residue
classes for their multiples and a tail decomposition, give m(n)
\textasciitilde{} (4/3)n\^{}(3/2)/log n.

\emph{The model's own caveat:} The argument is qualitative and does not
extract a useful rate of convergence.

\section{Model writeup}\label{model-writeup}

The text below is the model output verbatim, apart from moving its
machine-readable verdict block into the summary above and shifting
Markdown heading levels for this document. Mechanical TeX defects and
equation tags were normalized where needed for compilation.

\subsection{\texorpdfstring{A proof of the \(4/3\)
asymptotic}{A proof of the 4/3 asymptotic}}\label{a-proof-of-the-43-asymptotic}

All translations below are integer translations. I use interval
cardinality as the packing length; using interval diameter instead
changes the parameter by at most one.

Write \[
A_d(n)=\{d,2d,\ldots,\lfloor n/d\rfloor d\},
\qquad
x=\sqrt n,
\qquad
T_n=\frac{n^{3/2}}{\log n}.
\]

\subsubsection{Theorem}\label{theorem}

\[
m(n)=\left(\frac43+o(1)\right)T_n.
\]

The key new ingredient is that, for a fixed number of suitably selected
primes, all pairwise CRT intersection phases can simultaneously be
placed in prescribed arcs. Thus one can allow three or more convex hulls
to overlap, rather than impose the path condition from the previous
attempt.

The proof also handles the composite-indexed progressions explicitly.
The only analytic number-theoretic input is the prime number theorem.

\begin{center}\rule{0.5\linewidth}{0.5pt}\end{center}

\subsection{1. The lower bound}\label{the-lower-bound}

For completeness, I verify the prime obstruction used in the previous
attempt.

Let \(p,q\le x\) be distinct primes, and suppose \[
t+A_p(n),\qquad u+A_q(n)
\] are disjoint, with \(u\ge t\). Put \(r_p=\lfloor n/p\rfloor\). Then
\[
u-t>r_pp-pq.
\qquad\text{(1.1)}
\]

Indeed, if \(0\le \delta=u-t\le r_pp-pq\), choose \(j\in\{1,\ldots,p\}\)
such that \[
\delta+jq\equiv0\pmod p.
\] Then \[
i=\frac{\delta+jq}{p}
\] satisfies \(1\le i\le r_p\), while \(j\le p\le\lfloor n/q\rfloor\),
because \(pq\le n\). Hence \[
t+ip=u+jq,
\] a contradiction.

Now take a packing of all \(A_p(n)\), \(p\le x\), in an interval of
cardinality \(M\), and order their translation bases: \[
t_{p_1}\le\cdots\le t_{p_s}.
\] Summing (1.1) gives \[
t_{p_s}-t_{p_1}
\ge (s-1)n-\sum_{i=1}^{s-1}p_ip_{i+1}-O(n).
\] Here the rounding errors sum to \(O(n)\). Also \[
\sum_{i=1}^{s-1}p_ip_{i+1}
\le \frac12\sum_{i=1}^{s-1}(p_i^2+p_{i+1}^2)
\le\sum_{p\le x}p^2.
\] Translation bases of progressions contained in one interval of
cardinality \(M\) differ by at most \(M+x\). Consequently, \[
M\ge n\pi(x)-\sum_{p\le x}p^2-O(n).
\qquad\text{(1.2)}
\]

The prime number theorem and partial summation yield \[
\pi(x)\sim\frac{x}{\log x},
\qquad
\sum_{p\le x}p^2\sim\frac{x^3}{3\log x}.
\] Since \(\log n=2\log x\), (1.2) implies \[
m(n)\ge\left(\frac43-o(1)\right)T_n.
\qquad\text{(1.3)}
\]

The rest of the proof establishes the matching upper bound.

\begin{center}\rule{0.5\linewidth}{0.5pt}\end{center}

\subsection{2. Packing all multiples of one
modulus}\label{packing-all-multiples-of-one-modulus}

For an integer \(q\ge1\), call \[
\mathcal B_q(n)=\{A_{qa}(n):1\le a\le \lfloor n/q\rfloor\}
\] the bundle associated with \(q\).

We first give a deliberately crude but useful way to pack a bundle.

\subsubsection{Lemma 2.1}\label{lemma-2.1}

There is an absolute constant \(K_0\) such that, for every \(N\ge1\),
the family \[
\{A_a(N):1\le a\le N\}
\] can be partitioned into at most \[
R(N)\le K_0N^{3/5}
\qquad\text{(2.1)}
\] classes, each consisting of pairwise disjoint \textbf{unshifted}
progressions.

\paragraph{Proof}\label{proof}

Give each \(a\le\sqrt N\) its own class.

For \(a>\sqrt N\), form the intersection graph of the unshifted sets
\(A_a(N)\). A neighbor \(b\) of \(a\) divides some multiple \(ka\le N\).
Therefore its degree is at most \[
\sum_{k\le N/a}\tau(ka)
\le \sqrt N\max_{t\le N}\tau(t).
\] The elementary divisor bound \[
\tau(t)=O(t^{1/10})
\] then gives maximum degree \(O(N^{3/5})\), so greedy coloring proves
(2.1).

For clarity, the divisor bound follows directly from prime
factorization. For all sufficiently large primes \(p\), \[
e+1\le 2^e\le p^{e/10}.
\] The finitely many smaller primes contribute only an absolute
multiplicative constant, since \[
\sup_{e\ge0}\frac{e+1}{p^{e/10}}<\infty.
\] \(\square\)

Since \[
A_{qa}(n)=qA_a(\lfloor n/q\rfloor),
\qquad\text{(2.2)}
\] we obtain the following.

\subsubsection{Corollary 2.2}\label{corollary-2.2}

The bundle \(\mathcal B_q(n)\) can be packed in an interval of
cardinality at most \[
(n+q)\left\lceil\frac{R(\lfloor n/q\rfloor)}q\right\rceil.
\qquad\text{(2.3)}
\]

\paragraph{Proof}\label{proof-1}

Take the coloring from Lemma 2.1 and partition its colors into batches
of at most \(q\) colors.

In one batch, assign a different residue \(r\in\{0,\ldots,q-1\}\) to
each color, and translate every progression of that color by \(r\).
Different colors are disjoint modulo \(q\); progressions of the same
color remain disjoint by (2.2). Everything in the batch lies in
\([1,n+q]\).

Concatenate the batches. \(\square\)

For primes \(p\asymp\sqrt n\), (2.1) gives \[
R(\lfloor n/p\rfloor)=O(n^{3/10})=o(p).
\qquad\text{(2.4)}
\] The fact that only \(o(p)\) residue classes are needed will let an
entire bundle follow a single prime's CRT phase.

\begin{center}\rule{0.5\linewidth}{0.5pt}\end{center}

\subsection{3. Simultaneous CRT phase
control}\label{simultaneous-crt-phase-control}

Let \(\mathbb T=\mathbb R/\mathbb Z\), and write \(\|z\|_{\mathbb T}\)
for distance to the nearest integer.

For distinct primes \(p_i,p_j\) and residues \(r_i\bmod p_i\),
\(r_j\bmod p_j\), define \[
Z_{ij}(r_i,r_j)
=
\frac{\operatorname{CRT}_{p_i,p_j}(r_i,r_j)}{p_ip_j}
\pmod1,
\qquad\text{(3.1)}
\] where the CRT representative is chosen in \(\{0,\ldots,p_ip_j-1\}\).

Thus \(Z_{ij}\) specifies the locations, modulo \(p_ip_j\), at which the
two infinite residue classes intersect.

\subsubsection{Lemma 3.1 --- Generic prime
blocks}\label{lemma-3.1-generic-prime-blocks}

Fix \[
k\ge2,\qquad 0<a<b,\qquad 0<\eta<1.
\] There is a constant \(L=L(k,a,b,\eta)\) with the following property.

For all sufficiently large \(X\), every set of primes in \([aX,bX]\) can
be partitioned into \(k\)-element blocks and at most \(L\) exceptional
primes, so that each block \(p_1,\ldots,p_k\) has this property:

\begin{quote}
For arbitrary open arcs \(J_{ij}\subseteq\mathbb T\), \(1\le i<j\le k\),
each of length \(\eta\), there are residues \(r_i\bmod p_i\) satisfying
\[
Z_{ij}(r_i,r_j)\in J_{ij}
\qquad\text{for every }i<j.
\qquad\text{(3.2)}
\]
\end{quote}

\paragraph{Proof}\label{proof-2}

Put \(E=\binom{k}{2}\). For \(i\ne j\), let \[
\overline{p_j}^{(p_i)}
\] denote the inverse of \(p_j\) modulo \(p_i\). The CRT formula gives
\[
Z_{ij}(r_i,r_j)
=
\frac{\overline{p_j}^{(p_i)}r_i}{p_i}
+
\frac{\overline{p_i}^{(p_j)}r_j}{p_j}
\pmod1.
\qquad\text{(3.3)}
\]

Choose the \(r_i\) independently and uniformly. For an integer vector \[
h=(h_{ij})_{i<j}\in\mathbb Z^E,
\] writing \(h_{ji}=h_{ij}\), we have \[
\mathbb E\exp\!\left(2\pi i\sum_{i<j}h_{ij}Z_{ij}\right)
=
\prod_{i=1}^k
\mathbf 1\!\left[
\sum_{j\ne i}h_{ij}\overline{p_j}^{(p_i)}
\equiv0\pmod{p_i}
\right].
\qquad\text{(3.4)}
\]

\paragraph{A finite Fourier criterion}\label{a-finite-fourier-criterion}

There is an integer \(H=H(k,\eta)\) such that (3.2) holds whenever the
expectation in (3.4) vanishes for every nonzero \(h\) with \[
\|h\|_\infty\le H.
\qquad\text{(3.5)}
\]

To see this, choose a nonnegative smooth periodic function \(f\),
supported inside an arc of length \(\eta\), with \[
\mu=\int_{\mathbb T}f>0.
\] For prescribed arc centers \(c_{ij}\), consider \[
F(z)=\prod_{i<j}f(z_{ij}-c_{ij}).
\] Its constant Fourier coefficient is \(\mu^E\). Its Fourier series is
absolutely convergent, and the absolute sum of coefficients outside
\(\|h\|_\infty\le H\) tends to zero as \(H\to\infty\), uniformly in the
centers \(c_{ij}\).

Choose \(H\) so that this tail is less than \(\mu^E/2\). Under (3.5),
(3.4) then gives \[
\mathbb EF(Z)\ge\mu^E/2>0.
\] Some residue vector therefore satisfies (3.2).

\paragraph{There are few bad tuples}\label{there-are-few-bad-tuples}

Call an ordered tuple of distinct primes bad if some nonzero \(h\)
satisfying (3.5) makes all the congruences in (3.4) hold.

Fix such an \(h\), and choose \(i,j\) with \(h_{ij}\ne0\). Fix all
primes in the tuple except \(p_j\). The congruence at \(i\) has the form
\[
h_{ij}\overline{p_j}^{(p_i)}\equiv-c\pmod{p_i},
\qquad\text{(3.6)}
\] where \(c\) is fixed.

For \(p_i>H\), if \(c=0\), (3.6) has no solution. Otherwise it forces
\(p_j\) into a single residue class modulo \(p_i\). Since \[
p_i\ge aX,\qquad p_j\in[aX,bX],
\] there are only \(O_{a,b}(1)\) candidate integers for \(p_j\).

Consequently, in any set of \(M\) primes from this interval, the number
of bad ordered \(k\)-tuples is \[
O_{k,a,b,\eta}(M^{k-1}).
\] For \(M\) larger than a suitable constant, this is less than the
number \((M)_k\) of ordered distinct tuples. Hence a good block exists.

Repeatedly remove good blocks until only a bounded number of primes
remain. \(\square\)

A useful point is that this counting argument does \textbf{not} require
information about primes in arithmetic progressions with growing moduli.
Once the other primes are fixed, there are only boundedly many candidate
integers for the remaining prime.

\begin{center}\rule{0.5\linewidth}{0.5pt}\end{center}

\subsection{4. Keeping many residue classes close in CRT
phase}\label{keeping-many-residue-classes-close-in-crt-phase}

We need a simple simultaneous-approximation observation.

\subsubsection{Lemma 4.1 --- A large centered Bohr
set}\label{lemma-4.1-a-large-centered-bohr-set}

Fix \(d\ge1\) and \(\rho>0\). There is \(c=c(d,\rho)>0\) such that, for
every integer \(q\ge2\) and every \(c_1,\ldots,c_d\in\mathbb Z\), the
set \[
\left\{s\bmod q:
\left\|\frac{c_js}{q}\right\|_{\mathbb T}<\rho
\text{ for all }j
\right\}
\] has at least \(cq\) elements.

\paragraph{Proof}\label{proof-3}

Choose an integer \(M>1/\rho\). Partition \([0,1)^d\) into \(M^d\)
congruent boxes and consider the \(q\) points \[
\left(\frac{c_1s}{q},\ldots,\frac{c_ds}{q}\right)\pmod1,
\qquad s\bmod q.
\] One box contains at least \(q/M^d\) residues. Fix one of them,
\(s_0\), and subtract \(s_0\) from all the others. The resulting
distinct residues satisfy all the required inequalities. \(\square\)

Applied with \(q=p_i\) and coefficients \(\overline{p_j}^{(p_i)}\), this
gives linearly many residue perturbations whose effects on every CRT
phase involving \(p_i\) are small.

\begin{center}\rule{0.5\linewidth}{0.5pt}\end{center}

\subsection{5. Packing the main prime
range}\label{packing-the-main-prime-range}

Fix constants \[
0<\epsilon<1<C.
\] Let \(\mathcal G_{\epsilon,C}(n)\) be the family of all progressions
\(A_d(n)\) whose index has a prime divisor in \[
[\epsilon x,Cx].
\]

I will prove \[
\limsup_{n\to\infty}
\frac{m(\mathcal G_{\epsilon,C}(n))}{T_n}
\le
2\int_\epsilon^C(1-u^2)_+\,du.
\qquad\text{(5.1)}
\]

When bundles overlap, we may pack labelled copies first and then delete
duplicates. Thus it is enough to pack all the bundles
\(\mathcal B_p(n)\) in this prime range.

\subsubsection{5.1 Packing one good block}\label{packing-one-good-block}

Fix an integer \(k\ge2\) and \[
0<\nu<\epsilon^2/4.
\] Put \[
h(u)=(1-u^2)_+,
\qquad
\rho=\frac{\nu}{100C^2}.
\]

Consider a block of primes \[
p_1,\ldots,p_k\in[ux,vx]\subseteq[\epsilon x,Cx]
\] having the property in Lemma 3.1 for arcs of length \(\rho\).

Set \[
g=\left\lceil n(h(u)+\nu)\right\rceil,
\qquad
W=n+\lceil Cx\rceil+1,
\qquad
S_i=(i-1)g.
\] We will place the whole bundle associated with \(p_i\) inside \[
I_i=[S_i,S_i+W-1].
\]

For \(i<j\), if \(I_i\cap I_j\ne\varnothing\), its real length is at
most \(W-g\). For all sufficiently large \(n\), \begin{equation*}
\begin{aligned}
W-g
&\le \bigl(\min(u^2,1)-\nu/2\bigr)n\\
&\le p_ip_j-\nu n/2.
\end{aligned}
\qquad\text{(5.2)}
\end{equation*} Also \(p_ip_j\le C^2n\).

Project \(I_i\cap I_j\) into \(\mathbb T\) by dividing by \(p_ip_j\). By
(5.2), the complement of this forbidden arc has length at least \[
\frac{\nu}{2C^2}=50\rho.
\] Choose an open arc \(J_{ij}\) of length \(\rho\) in the complement,
at circular distance at least \(4\rho\) from the forbidden arc. If the
intervals are disjoint, choose any arc.

Lemma 3.1 supplies residues \(r_i\bmod p_i\) such that \[
Z_{ij}(r_i,r_j)\in J_{ij}
\qquad(i<j).
\qquad\text{(5.3)}
\]

Now put \[
N_i=\lfloor n/p_i\rfloor
\] and color the unshifted \(A_a(N_i)\) using Lemma 2.1. There are \[
R(N_i)=O_\epsilon(n^{3/10})=o(p_i)
\qquad\text{(5.4)}
\] colors.

For each \(i\), Lemma 4.1 supplies at least \(c_{k,\rho}p_i\) residues
\(s\bmod p_i\) satisfying \[
\left\|
\frac{\overline{p_j}^{(p_i)}s}{p_i}
\right\|_{\mathbb T}<\rho
\qquad(j\ne i).
\qquad\text{(5.5)}
\] By (5.4), there are enough distinct such residues to assign one,
denoted \(s_{i,c}\), to every color \(c\).

For color \(c\), choose \(t_{i,c}\in\{0,\ldots,p_i-1\}\) such that \[
S_i+t_{i,c}\equiv r_i+s_{i,c}\pmod{p_i}.
\] Translate every \(A_{ap_i}(n)\) of that color by \(S_i+t_{i,c}\).

These translates lie in \(I_i\). Within one bundle:

\begin{itemize}
\tightlist
\item
  progressions of the same color are disjoint by their unshifted
  coloring;
\item
  different colors occupy distinct residues modulo \(p_i\).
\end{itemize}

For two different bundles, their relevant CRT phase is \begin{equation*}
\begin{aligned}
Z_{ij}(r_i+s_{i,c},r_j+s_{j,c'})
={}&Z_{ij}(r_i,r_j)\\
&+\frac{\overline{p_j}^{(p_i)}s_{i,c}}{p_i}
+\frac{\overline{p_i}^{(p_j)}s_{j,c'}}{p_j}
\pmod1.
\end{aligned}
\qquad\text{(5.6)}
\end{equation*} By (5.5), this differs from the phase in (5.3) by
circular distance less than \(2\rho\). It therefore remains outside the
forbidden projection of \(I_i\cap I_j\).

A common point of two progressions would belong to both envelopes and
have precisely the CRT phase in (5.6), which is impossible.

We have proved that all bundles in this block can be packed in an
interval of cardinality at most \[
(k-1)g+W.
\qquad\text{(5.7)}
\]

\subsubsection{5.2 Summing the blocks}\label{summing-the-blocks}

Partition \([\epsilon,C]\) into finitely many intervals \[
\epsilon=u_0<u_1<\cdots<u_J=C.
\] For each bin, apply Lemma 3.1 to its primes, obtaining \(k\)-element
good blocks and \(O(1)\) exceptional primes.

An exceptional prime's whole bundle fits in an interval of cardinality
\(W\), since (5.4) eventually gives \(R(N_i)\le p_i\). Thus all
exceptions cost \(O(n)=o(T_n)\).

Let \[
M_j=\#\{p:u_{j-1}x\le p<u_jx\},
\qquad
g_j=\left\lceil n(h(u_{j-1})+\nu)\right\rceil.
\] Using (5.7) and concatenating blocks, \[
m(\mathcal G_{\epsilon,C}(n))
\le
\sum_{j=1}^J M_j\left(g_j+\frac Wk\right)+O(n).
\qquad\text{(5.8)}
\]

The prime number theorem gives \[
M_j=\left(u_j-u_{j-1}+o(1)\right)\frac{x}{\log x}.
\] Dividing (5.8) by \(T_n\), we get \[
\limsup_{n\to\infty}
\frac{m(\mathcal G_{\epsilon,C}(n))}{T_n}
\le
2\sum_{j=1}^J
(u_j-u_{j-1})
\left(h(u_{j-1})+\nu+\frac1k\right).
\]

All parameters here are fixed before taking \(n\to\infty\). We may
therefore let the partition mesh tend to zero, then let
\(\nu\downarrow0\) and \(k\to\infty\). This proves (5.1).

For \(C>1\), the resulting bound is \[
m(\mathcal G_{\epsilon,C}(n))
\le
\left(
\frac43-2\epsilon+\frac23\epsilon^3+o(1)
\right)T_n.
\qquad\text{(5.9)}
\]

\begin{center}\rule{0.5\linewidth}{0.5pt}\end{center}

\subsection{6. Two tools for the large-prime
tail}\label{two-tools-for-the-large-prime-tail}

\subsubsection{Lemma 6.1 --- Folding a scaled
packing}\label{lemma-6.1-folding-a-scaled-packing}

Suppose a family of finite integer sets, each of diameter at most \(N\),
has a packing in an interval of cardinality \(M\). For an integer
\(a\ge1\), the family obtained by multiplying every set by \(a\) has a
packing of cardinality at most \[
M+aN+a.
\qquad\text{(6.1)}
\]

\paragraph{Proof}\label{proof-4}

Translate the given packing into \([0,M-1]\), and put \[
H=\lceil M/a\rceil.
\] Assign a packed set \(B\) to the strip indexed by
\(j\in\{0,\ldots,a-1\}\) containing its minimum: \[
jH\le \min B<(j+1)H.
\] Replace it by \[
a(B-jH)+j.
\]

Sets from the same strip remain disjoint. Sets from different strips
have distinct residues modulo \(a\). All resulting sets lie in an
interval of cardinality at most \[
aH+aN\le M+aN+a.
\] \(\square\)

In particular, because \[
A_{ap}(n)=aA_p(\lfloor n/a\rfloor),
\qquad\text{(6.2)}
\] folding lets us scale a prime-indexed packing without multiplying its
leading packing length by \(a\).

\subsubsection{Lemma 6.2 --- Packing large prime
indices}\label{lemma-6.2-packing-large-prime-indices}

There is an absolute constant \(B\) such that, whenever
\(P\ge2\sqrt N\), \[
m_{\{p\text{ prime}:P<p\le N\}}(N)
\le
B\left(
\frac{N^2}{P\log P}+N\log(2N)
\right).
\qquad\text{(6.3)}
\] If \(P>N\), the family is empty.

\paragraph{Proof}\label{proof-5}

Use dyadic bins \((Q,2Q]\), starting with \(Q=P\). In one such bin, put
\[
K=\left\lfloor\frac{Q^2}{2N}\right\rfloor.
\] Since \(Q\ge2\sqrt N\), \[
K\ge\frac{Q^2}{4N}.
\qquad\text{(6.4)}
\]

Any group of at most \(K\) prime-indexed progressions from the bin can
be packed inside \([1,2N]\).

To see this, insert them successively. For a new prime \(p\), choose a
shift \(r\in\{0,\ldots,p-1\}\). Each previously occupied point forbids
at most one residue \(r\bmod p\). The previous progressions contain
fewer than \[
K\frac NQ\le \frac Q2<p
\] points, so some residue remains available. The resulting progression
lies in \([1,2N]\).

By (6.4), the total length for a bin is at most \[
2N\left(\frac{4N\pi(2Q)}{Q^2}+1\right).
\] Using \(\pi(t)=O(t/\log t)\) and summing over dyadic \(Q\), \[
\sum_{Q=P,2P,\ldots}\frac1{Q\log Q}
\le \frac2{P\log P}.
\] There are \(O(\log(2N))\) bins. These estimates give (6.3).
\(\square\)

\begin{center}\rule{0.5\linewidth}{0.5pt}\end{center}

\subsection{7. Covering every index}\label{covering-every-index}

Fix \(0<\epsilon<1\) and \(C>2\), and set \[
Y=\lceil n^{1/5}\rceil.
\] For large \(n\), \[
Y^2<\epsilon x.
\]

We cover all indices \(d\le n\) by four types of families:

\begin{enumerate}
\def\labelenumi{\arabic{enumi}.}
\tightlist
\item
  the main family \(\mathcal G_{\epsilon,C}(n)\);
\item
  bundles \(\mathcal B_q(n)\) for \[
  q\in\mathcal Q_n:=
  \{q\in\mathbb Z:Y\le q\le Y^2\}
  \cup
  \{p\text{ prime}:Y\le p<\epsilon x\};
  \]
\item
  individual progressions with \(d<Y\);
\item
  for each \(1\le a<Y\), the tail family \[
  \{A_{ap}(n):p\text{ prime},\ p>Cx,\ ap\le n\}.
  \]
\end{enumerate}

Here is a verification that this is a cover.

If \(d\) has a prime divisor in \([Y,Cx]\), it is covered by type 1 or
type 2.

Otherwise, every prime factor of \(d\) is either smaller than \(Y\) or
larger than \(Cx\). There is at most one factor larger than \(Cx\),
including multiplicity, because \(C>1\) and \(d\le n\). Hence \[
d=s\quad\text{or}\quad d=sp,
\] where all prime factors of \(s\) are smaller than \(Y\), and \(p>Cx\)
is prime when present.

If \(s\ge Y\), multiply its prime factors until the product first
reaches \(Y\). This gives a divisor \[
Y\le q<Y^2,
\] so type 2 applies. If \(s<Y\), then either \(d<Y\), giving type 3, or
\(d=sp\), giving type 4.

Thus no indices are omitted.

\subsubsection{7.1 Cost of the small-divisor
bundles}\label{cost-of-the-small-divisor-bundles}

By Corollary 2.2 and Lemma 2.1, \begin{equation*}
\begin{aligned}
M_{\mathcal Q}
&\le
(n+\epsilon x+1)
\sum_{q\in\mathcal Q_n}
\left(1+K_0n^{3/5}q^{-8/5}\right)\\
&\le
(n+\epsilon x+1)
\left(
\pi(\epsilon x)+Y^2+
O\!\left(n^{3/5}Y^{-3/5}\right)
\right).
\end{aligned}
\qquad\text{(7.1)}
\end{equation*} Now \[
Y^2=O(n^{2/5}),
\qquad
n^{3/5}Y^{-3/5}=O(n^{12/25}).
\] Both give \(o(T_n)\) after multiplication by \(n+O(x)\), since
\(37/25<3/2\). The prime number theorem therefore yields \[
M_{\mathcal Q}\le(2\epsilon+o(1))T_n.
\qquad\text{(7.2)}
\]

The individual progressions with \(d<Y\), packed separately, cost at
most \[
nY=o(T_n).
\qquad\text{(7.3)}
\]

\subsubsection{7.2 Cost of the large-prime
tail}\label{cost-of-the-large-prime-tail}

Put \[
P=Cx,\qquad N_a=\lfloor n/a\rfloor.
\] For every \(a<Y\), we have \(P\ge2\sqrt{N_a}\). By Lemma 6.2,
followed by folding and (6.2), the \(a\)-th tail family can be packed in
length at most \[
B\left(
\frac{N_a^2}{P\log P}+N_a\log(2N_a)
\right)+n+a.
\] Summing over \(a<Y\), and using \[
\sum_{a\ge1}a^{-2}<2,
\qquad
\sum_{a<Y}a^{-1}=O(\log Y),
\] gives \[
M_{\mathrm{tail}}
=
O\!\left(\frac{n^2}{P\log P}\right)
+
O(n\log^2 n+nY+Y^2).
\qquad\text{(7.4)}
\] The second term is \(o(T_n)\). Since \[
\frac{n^2}{Cx\log(Cx)}
\le \frac2C\,T_n,
\] there is an absolute constant \(B_0\), independent of \(\epsilon,C\),
such that \[
M_{\mathrm{tail}}
\le
\left(\frac{B_0}{C}+o(1)\right)T_n.
\qquad\text{(7.5)}
\]

\begin{center}\rule{0.5\linewidth}{0.5pt}\end{center}

\subsection{8. Completing the upper
bound}\label{completing-the-upper-bound}

Concatenate the packings of the covering families, then delete any
duplicate copies.

Combining (5.9), (7.2), (7.3), and (7.5), \begin{equation*}
\begin{aligned}
\limsup_{n\to\infty}\frac{m(n)}{T_n}
&\le
\left(\frac43-2\epsilon+\frac23\epsilon^3\right)
+2\epsilon+\frac{B_0}{C}\\
&=
\frac43+\frac23\epsilon^3+\frac{B_0}{C}.
\end{aligned}
\end{equation*} This holds for every fixed \(0<\epsilon<1\) and \(C>2\).
Letting \(\epsilon\downarrow0\) and \(C\to\infty\), \[
\limsup_{n\to\infty}\frac{m(n)}{T_n}\le\frac43.
\] Together with (1.3), \[
\boxed{\displaystyle
m(n)\sim\frac43\,\frac{n^{3/2}}{\log n}.}
\]

\subsubsection{What closes the earlier
gaps}\label{what-closes-the-earlier-gaps}

Two points distinguish this argument from the previous path-based
attempt:

\begin{itemize}
\tightlist
\item
  \textbf{Multiple overlapping hulls:} Lemma 3.1 simultaneously controls
  every pairwise CRT phase in a fixed block. No triple-overlap
  restriction is imposed.
\item
  \textbf{Composite indices:} Lemma 2.1 places an entire prime bundle in
  \(o(p)\) residue classes. Lemma 4.1 lets all these classes remain
  close to the bundle's chosen CRT phase. The indices not captured by
  main-range bundles are explicitly covered and bounded in Section 7.
\end{itemize}

All block sizes and phase tolerances are fixed before \(n\to\infty\).
This order of limits is sufficient for the asymptotic, although it
produces no practical convergence rate.

\appendix

\section{Provenance and review
status}\label{provenance-and-review-status}

\begin{warning}

\textbf{UNREFEREED MODEL OUTPUT.} This document was selected solely
because GPT-6 Astra labelled its own result
\texttt{would\_publish:\ true} and returned \texttt{proved} or
\texttt{disproved}. It has not passed the adversarial LLM referee used
for the earlier Sol campaign, has not been checked by a human
mathematician, and has not been checked for novelty. Treat every
mathematical and bibliographic claim as unverified.

\end{warning}

{\def\LTcaptype{none} % do not increment counter
\begin{longtable}[]{@{}
  >{\raggedright\arraybackslash}p{(\linewidth - 2\tabcolsep) * \real{0.5000}}
  >{\raggedright\arraybackslash}p{(\linewidth - 2\tabcolsep) * \real{0.5000}}@{}}
\toprule\noalign{}
\endhead
\bottomrule\noalign{}
\endlastfoot
Catalog id & \texttt{2603.02786\_\_00} \\
Catalog entry &
\href{https://graph-theory-ai.github.io/graph-conjectures/arxiv/2603.02786__00/}{Conjecture
1} \\
Source corpus & arxiv \\
Campaign leg & selected second attempts (\texttt{attacks\_retry}) \\
Source paper / entry & Packing arithmetic progressions \\
Model verdict & \textbf{proved} (confidence: high) \\
Selection criterion & \texttt{would\_publish:\ true} and verdict
\texttt{proved} or \texttt{disproved} \\
Model & \texttt{gpt-6-astra}, reasoning effort \texttt{max},
\texttt{mode=pro}, flex service tier \\
Original artifact &
\texttt{attacks\_retry/2603.02786\_\_00/output.md} \\
Independent review & \textbf{None} \\
arXiv & \href{https://arxiv.org/abs/2603.02786}{arXiv:2603.02786} \\
Previous attempt & \texttt{attacks/2603.02786\_\_00}: gpt-5.6-sol
verdict \texttt{partial} \\
Model's caveats & The argument is qualitative and does not extract a
useful rate of convergence. \\
\end{longtable}
}

The generated Markdown and LaTeX sources for this PDF are stored under
\texttt{to\_review\_astra/src/2603.02786\_\_00/}.

\end{document}
